% GERADO POR lab/sumario.py — NÃO EDITE À MÃO.
% Uma cópia do livro.toc da compilação anterior: o \tableofcontents não pode ser
% repetido no fim do documento, porque o arquivo ainda está sendo escrito.
\contentsline {chapter}{Como este livro mede}{11}{chapter*.2}%
\contentsline {chapter}{Introdução}{13}{chapter*.3}%
\contentsline {part}{I\hspace {1em}A raiz}{15}{part.1}%
\contentsline {chapter}{\numberline {1}O número e a barra}{19}{chapter.1}%
\contentsline {section}{\numberline {1.1}Dois números para a mesma pergunta}{19}{section.1.1}%
\contentsline {section}{\numberline {1.2}O sorteio, a lei e a variável}{20}{section.1.2}%
\contentsline {section}{\numberline {1.3}A média de muitos sorteios}{21}{section.1.3}%
\contentsline {section}{\numberline {1.4}A barra que se confere}{22}{section.1.4}%
\contentsline {section}{\numberline {1.5}A barra no dado, e o que ela deixa de fora}{24}{section.1.5}%
\contentsline {section}{\numberline {1.6}O que fica}{25}{section.1.6}%
\contentsline {chapter}{\numberline {2}O que a soma guarda}{27}{chapter.2}%
\contentsline {section}{\numberline {2.1}A conta que fecha no mês e mente em vinte e seis anos}{27}{section.2.1}%
\contentsline {section}{\numberline {2.2}O nível e o incremento}{28}{section.2.2}%
\contentsline {section}{\numberline {2.3}O logaritmo, e por que o livro mede nele}{28}{section.2.3}%
\contentsline {section}{\numberline {2.4}O desvio do nível cresce com a raiz do horizonte}{29}{section.2.4}%
\contentsline {section}{\numberline {2.5}A mesma previsão na série real}{30}{section.2.5}%
\contentsline {section}{\numberline {2.6}O que fica}{32}{section.2.6}%
\contentsline {chapter}{\numberline {3}A medida, o corte e o preço}{35}{chapter.3}%
\contentsline {section}{\numberline {3.1}Quanto eu posso perder amanhã}{35}{section.3.1}%
\contentsline {section}{\numberline {3.2}O dado e a regra}{36}{section.3.2}%
\contentsline {section}{\numberline {3.3}A conta que o corte faz sozinho}{37}{section.3.3}%
\contentsline {section}{\numberline {3.4}Os blocos têm data}{39}{section.3.4}%
\contentsline {section}{\numberline {3.5}O preço muda de lugar}{41}{section.3.5}%
\contentsline {section}{\numberline {3.6}Uma linha erguida no passado}{42}{section.3.6}%
\contentsline {section}{\numberline {3.7}O que fica}{43}{section.3.7}%
\contentsline {part}{II\hspace {1em}As voltas}{45}{part.2}%
\contentsline {chapter}{\numberline {4}O orçamento do alarme}{49}{chapter.4}%
\contentsline {section}{\numberline {4.1}O vigia}{49}{section.4.1}%
\contentsline {section}{\numberline {4.2}Quanto ele soa quando nada muda}{50}{section.4.2}%
\contentsline {section}{\numberline {4.3}Quanto ele demora quando algo muda}{52}{section.4.3}%
\contentsline {section}{\numberline {4.4}A troca}{53}{section.4.4}%
\contentsline {section}{\numberline {4.5}O piso}{54}{section.4.5}%
\contentsline {section}{\numberline {4.6}O mesmo corte, lido mais fundo}{55}{section.4.6}%
\contentsline {section}{\numberline {4.7}O que fica}{58}{section.4.7}%
\contentsline {chapter}{\numberline {5}A forma decide o preço}{61}{chapter.5}%
\contentsline {section}{\numberline {5.1}Três formas de mudança}{61}{section.5.1}%
\contentsline {section}{\numberline {5.2}O que a memória faz}{63}{section.5.2}%
\contentsline {section}{\numberline {5.3}O que fica}{67}{section.5.3}%
\contentsline {chapter}{\numberline {6}O dia da semana}{69}{chapter.6}%
\contentsline {section}{\numberline {6.1}O domingo não é uma mudança}{69}{section.6.1}%
\contentsline {section}{\numberline {6.2}O tamanho do calendário}{70}{section.6.2}%
\contentsline {section}{\numberline {6.3}Ler o dia no seu próprio dia}{71}{section.6.3}%
\contentsline {section}{\numberline {6.4}A partição divide a memória}{71}{section.6.4}%
\contentsline {section}{\numberline {6.5}O que a memória compra de volta}{73}{section.6.5}%
\contentsline {section}{\numberline {6.6}Onde não há calendário, a partição só custa}{75}{section.6.6}%
\contentsline {section}{\numberline {6.7}O que fica}{76}{section.6.7}%
\contentsline {chapter}{\numberline {7}O dia em que tudo cai junto}{77}{chapter.7}%
\contentsline {section}{\numberline {7.1}O que a margem não vê}{77}{section.7.1}%
\contentsline {section}{\numberline {7.2}A conta que a soma faz}{78}{section.7.2}%
\contentsline {section}{\numberline {7.3}O par é que decide}{81}{section.7.3}%
\contentsline {section}{\numberline {7.4}O controle}{82}{section.7.4}%
\contentsline {section}{\numberline {7.5}O que a proteção individual não cobre}{82}{section.7.5}%
\contentsline {section}{\numberline {7.6}O que fica}{83}{section.7.6}%
\contentsline {chapter}{\numberline {8}A estatística que não anuncia}{85}{chapter.8}%
\contentsline {section}{\numberline {8.1}A conta que todo mundo faz}{85}{section.8.1}%
\contentsline {section}{\numberline {8.2}Três sinais e o que vem depois}{86}{section.8.2}%
\contentsline {section}{\numberline {8.3}O controle que desfaz a leitura}{87}{section.8.3}%
\contentsline {section}{\numberline {8.4}O que sobra}{89}{section.8.4}%
\contentsline {section}{\numberline {8.5}O que fica}{90}{section.8.5}%
\contentsline {chapter}{\numberline {9}Quantas explicações cabem}{93}{chapter.9}%
\contentsline {section}{\numberline {9.1}A pergunta que o capítulo anterior deixou}{93}{section.9.1}%
\contentsline {section}{\numberline {9.2}A família mais simples que existe}{94}{section.9.2}%
\contentsline {section}{\numberline {9.3}Com a taxa sozinha, quase tudo cabe}{95}{section.9.3}%
\contentsline {section}{\numberline {9.4}O pior bloco esvazia o conjunto}{95}{section.9.4}%
\contentsline {section}{\numberline {9.5}O estado que dura}{96}{section.9.5}%
\contentsline {section}{\numberline {9.6}O tamanho da região é de quem mede}{98}{section.9.6}%
\contentsline {section}{\numberline {9.7}O que fica}{99}{section.9.7}%
\contentsline {chapter}{\numberline {10}Ver mundos basta}{101}{chapter.10}%
\contentsline {section}{\numberline {10.1}A tentativa de sempre}{101}{section.10.1}%
\contentsline {section}{\numberline {10.2}O que os três concordam}{102}{section.10.2}%
\contentsline {section}{\numberline {10.3}O que os três discordam}{102}{section.10.3}%
\contentsline {section}{\numberline {10.4}Com a lei compartilhada}{104}{section.10.4}%
\contentsline {section}{\numberline {10.5}Qual estatística é o gargalo}{104}{section.10.5}%
\contentsline {section}{\numberline {10.6}A tentativa de consertar, e o que ela ensinou}{105}{section.10.6}%
\contentsline {section}{\numberline {10.7}O que fica}{106}{section.10.7}%
\contentsline {part}{III\hspace {1em}A travessia}{109}{part.3}%
\contentsline {chapter}{\numberline {11}O vigia não vê o que não está na série}{113}{chapter.11}%
\contentsline {section}{\numberline {11.1}A primeira travessia}{113}{section.11.1}%
\contentsline {section}{\numberline {11.2}A construção}{114}{section.11.2}%
\contentsline {section}{\numberline {11.3}A margem é cega, e não por calibração}{114}{section.11.3}%
\contentsline {section}{\numberline {11.4}O instrumento que vê, e o que ele cobra}{115}{section.11.4}%
\contentsline {section}{\numberline {11.5}A comparação que dá sentido ao número}{116}{section.11.5}%
\contentsline {section}{\numberline {11.6}O controle}{117}{section.11.6}%
\contentsline {section}{\numberline {11.7}O que fica}{118}{section.11.7}%
\contentsline {chapter}{\numberline {12}O desenho da intervenção}{119}{chapter.12}%
\contentsline {section}{\numberline {12.1}O experimento que não separa}{119}{section.12.1}%
\contentsline {section}{\numberline {12.2}O ato e a resposta}{120}{section.12.2}%
\contentsline {section}{\numberline {12.3}A medição}{123}{section.12.3}%
\contentsline {section}{\numberline {12.4}O preço de separar}{125}{section.12.4}%
\contentsline {section}{\numberline {12.5}A promessa do desenho, medida}{127}{section.12.5}%
\contentsline {section}{\numberline {12.6}O mesmo orçamento no lugar errado}{128}{section.12.6}%
\contentsline {section}{\numberline {12.7}O que fica}{128}{section.12.7}%
\contentsline {chapter}{\numberline {13}A partilha do relógio}{131}{chapter.13}%
\contentsline {section}{\numberline {13.1}A barreira tem dois ingredientes}{131}{section.13.1}%
\contentsline {section}{\numberline {13.2}As duas pressas}{132}{section.13.2}%
\contentsline {section}{\numberline {13.3}A partilha}{134}{section.13.3}%
\contentsline {section}{\numberline {13.4}Quando a partilha morde}{135}{section.13.4}%
\contentsline {section}{\numberline {13.5}O que fica}{136}{section.13.5}%
\contentsline {chapter}{\numberline {14}A memória que aprende}{139}{chapter.14}%
\contentsline {section}{\numberline {14.1}A janela de um ano, num mundo que muda}{139}{section.14.1}%
\contentsline {section}{\numberline {14.2}Esquecer tem taxa}{140}{section.14.2}%
\contentsline {section}{\numberline {14.3}A medição}{142}{section.14.3}%
\contentsline {section}{\numberline {14.4}O que cabe prometer}{143}{section.14.4}%
\contentsline {section}{\numberline {14.5}O que fica}{145}{section.14.5}%
\contentsline {chapter}{\numberline {15}O vigia da direção}{147}{chapter.15}%
\contentsline {section}{\numberline {15.1}O orçamento que a conta promete}{147}{section.15.1}%
\contentsline {section}{\numberline {15.2}O instrumento, conferido}{149}{section.15.2}%
\contentsline {section}{\numberline {15.3}Quanto tempo para ver}{149}{section.15.3}%
\contentsline {section}{\numberline {15.4}O ano não é a direção}{151}{section.15.4}%
\contentsline {section}{\numberline {15.5}O que fica}{152}{section.15.5}%
\contentsline {chapter}{\numberline {16}Quem cai primeiro}{155}{chapter.16}%
\contentsline {section}{\numberline {16.1}A ordem que a proteção supõe}{155}{section.16.1}%
\contentsline {section}{\numberline {16.2}A contagem}{156}{section.16.2}%
\contentsline {section}{\numberline {16.3}O nulo, e o par}{157}{section.16.3}%
\contentsline {section}{\numberline {16.4}Quanto dura o adiantamento}{157}{section.16.4}%
\contentsline {section}{\numberline {16.5}Se repete}{159}{section.16.5}%
\contentsline {section}{\numberline {16.6}O que fica}{160}{section.16.6}%
\contentsline {part}{IV\hspace {1em}O que não se compra}{163}{part.4}%
\contentsline {chapter}{\numberline {17}O mundo que faltou}{167}{chapter.17}%
\contentsline {section}{\numberline {17.1}O pior dia como teto}{167}{section.17.1}%
\contentsline {section}{\numberline {17.2}A conta exata}{168}{section.17.2}%
\contentsline {section}{\numberline {17.3}A conta conferida, e o dado real}{169}{section.17.3}%
\contentsline {section}{\numberline {17.4}O que a mudança faz com a conta}{170}{section.17.4}%
\contentsline {section}{\numberline {17.5}O que fica}{172}{section.17.5}%
\contentsline {chapter}{\numberline {18}O que foi apagado}{173}{chapter.18}%
\contentsline {section}{\numberline {18.1}A tentativa que explode}{173}{section.18.1}%
\contentsline {section}{\numberline {18.2}As duas reconstruções}{174}{section.18.2}%
\contentsline {section}{\numberline {18.3}O horizonte, e o que ele não é}{176}{section.18.3}%
\contentsline {section}{\numberline {18.4}Bits não compram o passado}{177}{section.18.4}%
\contentsline {section}{\numberline {18.5}O que fica}{178}{section.18.5}%
\contentsline {chapter}{\numberline {19}O que a margem deixa de fora}{181}{chapter.19}%
\contentsline {section}{\numberline {19.1}O experimento controlado pelo segundo momento}{181}{section.19.1}%
\contentsline {section}{\numberline {19.2}O que fica parado e o que cresce}{182}{section.19.2}%
\contentsline {section}{\numberline {19.3}A margem cumpre a taxa e paga mais fundo}{184}{section.19.3}%
\contentsline {section}{\numberline {19.4}O que fica}{186}{section.19.4}%
\contentsline {part}{V\hspace {1em}O enquadramento}{187}{part.5}%
\contentsline {chapter}{\numberline {20}O pedaço escolhido}{191}{chapter.20}%
\contentsline {section}{\numberline {20.1}A resposta de um pedaço só}{191}{section.20.1}%
\contentsline {section}{\numberline {20.2}Nem tudo balança}{193}{section.20.2}%
\contentsline {section}{\numberline {20.3}O que fica}{195}{section.20.3}%
\contentsline {chapter}{\numberline {21}O agora}{197}{chapter.21}%
\contentsline {section}{\numberline {21.1}O dado que chega velho}{197}{section.21.1}%
\contentsline {section}{\numberline {21.2}O atraso não custa dias}{198}{section.21.2}%
\contentsline {section}{\numberline {21.3}O que ele cobra, então}{199}{section.21.3}%
\contentsline {section}{\numberline {21.4}O agora é uma janela}{200}{section.21.4}%
\contentsline {section}{\numberline {21.5}O que fica}{202}{section.21.5}%
\contentsline {part}{VI\hspace {1em}O fecho}{203}{part.6}%
\contentsline {chapter}{\numberline {22}O fecho: a promessa com os buracos abertos}{207}{chapter.22}%
\contentsline {section}{\numberline {22.1}A medição da pilha}{208}{section.22.1}%
\contentsline {section}{\numberline {22.2}Os buracos se compõem}{209}{section.22.2}%
\contentsline {section}{\numberline {22.3}O balanço}{210}{section.22.3}%
\contentsline {section}{\numberline {22.4}O que fica}{212}{section.22.4}%
\contentsline {chapter}{Conclusão}{215}{chapter*.4}%
\contentsline {chapter}{Referências}{217}{chapter*.5}%
\contentsline {chapter}{Sumário completo}{249}{chapter*.6}%
\contentsline {chapter}{Lista de figuras}{257}{chapter*.6}%
\contentsline {chapter}{Lista de tabelas}{269}{chapter*.7}%
\contentsline {chapter}{Proposições e definições}{275}{chapter*.8}%
